\documentclass[12pt]{article}
\usepackage[utf8]{inputenc}
\usepackage[T1]{fontenc}
\usepackage{hyperref}
\title{Fon Language Resources for K-12 Curriculum Development}
\author{Compiled by Claude Sonnet 5.5 (claude-sonnet-5-5) for Andr\'as Kornai}
\date{}
\begin{document}
\maketitle

\section{Introduction}
Fon (Fongbe, ISO 639-3 \texttt{fon}) is a Gbe language spoken mainly in Benin. This report covers only resources that appeared in the web searches. Few school-level materials were found.

\section{Primary Educational Resources}
No Fon primer, textbook or official school alphabet document was found in the results. Wikipedia notes that French was long the sole language of instruction in Benin and that some subjects have been piloted in local languages including Fon: \url{https://en.wikipedia.org/wiki/Fon_language}.

\section{Linguistic and Reference Materials}
\begin{itemize}
\item Koeppe: \emph{Dictionnaire Fon--Fran\c{c}ais} (2003), about 8,000 entries with a condensed grammar of about fifty pages, using the standard Fon orthography with tone marking: \url{https://www.koeppe.de/titel_dictionnaire-fon-francais}.
\item Koeppe: \emph{Dictionnaire Fon--Fran\c{c}ais / Fran\c{c}ais--Fon}, the combined listing: \url{https://www.koeppe.de/titel_print_dictionnaire-fon-francais-francais-fon}.
\item Overview of the language, phonology and writing: \url{https://en.wikipedia.org/wiki/Fon_language}.
\end{itemize}

\section{Religious and Community Resources}
None found in the results.

\section{Digital and Audio Resources}
None found in the results.

\section{Curriculum Development Context}
Fon is a language of wider communication, and a standard orthography exists. Searches did not surface the Beninese national-language curriculum or Fon textbooks; those, along with the national orthography guide, should be sought through Benin education authorities.
\end{document}